% 附录A：常用公式表

\chapter{常用公式表}

\section{欧拉公式}
\begin{zhongyao}[欧拉公式]
\begin{align}
    \mathrm{e}^{\mathrm{i}\theta} &= \cos\theta + \mathrm{i}\sin\theta \\
    \cos\theta &= \frac{\mathrm{e}^{\mathrm{i}\theta} + \mathrm{e}^{-\mathrm{i}\theta}}{2} \\
    \sin\theta &= \frac{\mathrm{e}^{\mathrm{i}\theta} - \mathrm{e}^{-\mathrm{i}\theta}}{2\mathrm{i}}
\end{align}
\end{zhongyao}

\section{三角恒等式}

\subsection{和差化积}
\begin{align}
    \sin\alpha + \sin\beta &= 2\sin\frac{\alpha+\beta}{2}\cos\frac{\alpha-\beta}{2} \\
    \sin\alpha - \sin\beta &= 2\cos\frac{\alpha+\beta}{2}\sin\frac{\alpha-\beta}{2} \\
    \cos\alpha + \cos\beta &= 2\cos\frac{\alpha+\beta}{2}\cos\frac{\alpha-\beta}{2} \\
    \cos\alpha - \cos\beta &= -2\sin\frac{\alpha+\beta}{2}\sin\frac{\alpha-\beta}{2}
\end{align}

\subsection{积化和差}
\begin{align}
    \sin\alpha\cos\beta &= \frac{1}{2}[\sin(\alpha+\beta) + \sin(\alpha-\beta)] \\
    \cos\alpha\sin\beta &= \frac{1}{2}[\sin(\alpha+\beta) - \sin(\alpha-\beta)] \\
    \cos\alpha\cos\beta &= \frac{1}{2}[\cos(\alpha+\beta) + \cos(\alpha-\beta)] \\
    \sin\alpha\sin\beta &= -\frac{1}{2}[\cos(\alpha+\beta) - \cos(\alpha-\beta)]
\end{align}

\subsection{倍角与半角}
\begin{align}
    \sin 2\theta &= 2\sin\theta\cos\theta \\
    \cos 2\theta &= \cos^2\theta - \sin^2\theta = 2\cos^2\theta - 1 = 1 - 2\sin^2\theta \\
    \sin 3\theta &= 3\sin\theta - 4\sin^3\theta \\
    \cos 3\theta &= 4\cos^3\theta - 3\cos\theta
\end{align}

\section{双曲函数}
\begin{align}
    \sinh x &= \frac{\mathrm{e}^x - \mathrm{e}^{-x}}{2} \\
    \cosh x &= \frac{\mathrm{e}^x + \mathrm{e}^{-x}}{2} \\
    \tanh x &= \frac{\sinh x}{\cosh x} = \frac{\mathrm{e}^x - \mathrm{e}^{-x}}{\mathrm{e}^x + \mathrm{e}^{-x}}
\end{align}

\textbf{恒等式}：
\begin{align}
    \cosh^2 x - \sinh^2 x &= 1 \\
    \sinh(x+y) &= \sinh x\cosh y + \cosh x\sinh y \\
    \cosh(x+y) &= \cosh x\cosh y + \sinh x\sinh y
\end{align}

\section{柯西-黎曼条件}
\begin{zhongyao}[柯西-黎曼条件]
\[
    \frac{\partial u}{\partial x} = \frac{\partial v}{\partial y}, \quad \frac{\partial u}{\partial y} = -\frac{\partial v}{\partial x}
\]
\end{zhongyao}

\section{泰勒展开}
\begin{align}
    \mathrm{e}^x &= 1 + x + \frac{x^2}{2!} + \frac{x^3}{3!} + \cdots \\
    \sin x &= x - \frac{x^3}{3!} + \frac{x^5}{5!} - \cdots \\
    \cos x &= 1 - \frac{x^2}{2!} + \frac{x^4}{4!} - \cdots \\
    \ln(1+x) &= x - \frac{x^2}{2} + \frac{x^3}{3} - \cdots, \quad |x| < 1 \\
    (1+x)^\alpha &= 1 + \alpha x + \frac{\alpha(\alpha-1)}{2!}x^2 + \cdots, \quad |x| < 1
\end{align}

\section{常用积分}
\begin{align}
    \oint_{|z|=R} \frac{\mathrm{d}z}{z - a} &= \begin{cases} 2\pi\mathrm{i}, & |a| < R \\ 0, & |a| > R \end{cases} \\
    \int_0^{+\infty} \frac{\sin x}{x}\,\mathrm{d}x &= \frac{\pi}{2} \\
    \int_{-\infty}^{+\infty} \mathrm{e}^{-x^2}\,\mathrm{d}x &= \sqrt{\pi} \\
    \int_0^{+\infty} \mathrm{e}^{-ax}\,\mathrm{d}x &= \frac{1}{a}, \quad a > 0 \\
    \int_0^{+\infty} x^n\mathrm{e}^{-ax}\,\mathrm{d}x &= \frac{n!}{a^{n+1}}, \quad a > 0
\end{align}

\section{积分公式}

\subsection{格林公式}
\begin{zhongyao}[格林公式]
设 $D$ 是平面上由分段光滑曲线 $L$ 围成的有界闭区域，$P(x,y)$、$Q(x,y)$ 在 $D$ 上有连续的一阶偏导数，则
\[
    \oint_L P\,\mathrm{d}x + Q\,\mathrm{d}y = \iint_D \left(\frac{\partial Q}{\partial x} - \frac{\partial P}{\partial y}\right)\mathrm{d}x\,\mathrm{d}y
\]
\end{zhongyao}

\subsection{高斯公式}
\begin{zhongyao}[高斯公式]
设 $V$ 是由分片光滑闭曲面 $S$ 围成的有界闭区域，$\boldsymbol{F}$ 有连续的一阶偏导数，则
\[
    \oint_S \boldsymbol{F} \cdot \mathrm{d}\boldsymbol{S} = \iiint_V \nabla \cdot \boldsymbol{F} \, \mathrm{d}V
\]
\end{zhongyao}

\subsection{斯托克斯公式}
\begin{zhongyao}[斯托克斯公式]
设 $S$ 是分片光滑的有向曲面，$L$ 是 $S$ 的边界曲线，则
\[
    \oint_L \boldsymbol{F} \cdot \mathrm{d}\boldsymbol{l} = \iint_S (\nabla \times \boldsymbol{F}) \cdot \mathrm{d}\boldsymbol{S}
\]
\end{zhongyao}

\section{三类方程的典型解}

\textbf{波动方程}（达朗贝尔公式）：
\begin{zhongyao}[达朗贝尔公式]
\[
    u(x, t) = \frac{1}{2}[\varphi(x-at) + \varphi(x+at)] + \frac{1}{2a}\int_{x-at}^{x+at} \psi(s)\,\mathrm{d}s
\]
\end{zhongyao}

\textbf{热传导方程}（泊松公式）：
\begin{equation}
    u(x, t) = \frac{1}{2a\sqrt{\pi t}}\int_{-\infty}^{+\infty} \varphi(\xi)\mathrm{e}^{-(x-\xi)^2/(4a^2t)}\,\mathrm{d}\xi
\end{equation}

\textbf{拉普拉斯方程}（基本解）：
\begin{equation}
    u = \frac{1}{r} \quad \text{（三维）}, \qquad u = \ln\frac{1}{r} \quad \text{（二维）}
\end{equation}
